% Preamble of the PDF vignettes (common layer; synced). gtm.sty, the author's manuscript layout (footer with date and
% "Page n of N", running head, full-width small captions), is used when LaTeX can find it: installed in TEXMFHOME,
% or GTM_STY set for tools/render-pdf.R. It is never copied into the package. Without it: the plain article class.
\IfFileExists{gtm.sty}{\usepackage{gtm}}{}
\providecommand{\runningtitle}{}
\renewcommand{\runningtitle}{DiallelAnalysisR}
% XeLaTeX loads unicode-math, where amsmath's \boldsymbol fails on Greek letters: use \symbf with the TeX bold
% style (bold italic lowercase Greek, bold upright capitals), as in the HTML
\makeatletter
\AtBeginDocument{\@ifpackageloaded{unicode-math}{\unimathsetup{bold-style=TeX}\renewcommand{\boldsymbol}[1]{\symbf{#1}}}{}}
\makeatother
% code and printed output (one Highlighting block per chunk, collapse = TRUE): long lines break instead of running
% into the margin (in HTML they scroll); \small fits the 100-character lines of the code style
\usepackage{fvextra}
\RecustomVerbatimEnvironment{Highlighting}{Verbatim}{breaklines,breakanywhere,commandchars=\\\{\},fontsize=\small}
% the vignette macros win over a command of the same name from another package (e.g. \Call of algpseudocode)
\newcommand{\vmacro}[1]{\providecommand{#1}{}\renewcommand{#1}}
% generated from _DiallelAnalysisR-macros.txt by tools/render-pdf.R: do not edit
\vmacro{\mY}{\mathbf{Y}}
\vmacro{\mT}{\mathbf{T}}
\vmacro{\mS}{\mathbf{S}}
\vmacro{\mR}{\mathbf{R}}
\vmacro{\vg}{\mathbf{g}}
\vmacro{\E}{\mathrm{E}}
\vmacro{\Var}{\mathrm{Var}}
\vmacro{\Cov}{\mathrm{Cov}}

